\documentclass[12pt]{article}


\usepackage[dvips,letterpaper,margin=0.75in,bottom=0.5in]{geometry}
\usepackage{cite}
\usepackage{slashed}
\usepackage{graphicx}
\usepackage{amsmath}
\usepackage{braket}
\usepackage{feynman}

\DeclareMathOperator{\Tr}{Tr}
\newcommand{\gv} {\ensuremath{g_{\mathrm v}}}
\newcommand{\ga} {\ensuremath{g_{\mathrm a}}}
\newcommand{\Ps} {\ensuremath{\slashed{P}}}
\newcommand{\Qs} {\ensuremath{\slashed{Q}}}
\newcommand{\Rs} {\ensuremath{\slashed{R}}}
\newcommand{\Ss} {\ensuremath{\slashed{S}}}
\newcommand{\ps}{\ensuremath{\slashed{p}}}
\newcommand{\GeV} {\ensuremath{\mathrm{Ge\kern -0.08em V}}}

\begin{document}

\title{Homework Assignment 7 \\ QED }

\maketitle

\noindent
All problems are graded on effort only.\\

\noindent
{\bf Problem 1:}  For the plane wave solutions to the Dirac Equation:
$$\def\arraystretch{1.7}
u^{(1)} = N \begin{pmatrix} 1 \\ 0 \\ {\displaystyle \frac{c p_z}{E+mc^2}} \\[5pt] {\displaystyle \frac{c p_x + i c p_y}{E + mc^2}} \end{pmatrix}, 
\hspace{1cm}
u^{(2)} = N \begin{pmatrix} 0 \\ 1 \\ {\displaystyle  \frac{c p_x - i c p_y}{E+mc^2}} \\ {\displaystyle  \frac{-c p_z}{E+mc^2}} \end{pmatrix}.
$$
Show that the normalization condition:
$$u^{\dagger} u = 2E/c$$
requires:
$$|N|^2 = (E + mc^2)/c$$
we choose to set:
$$N = \sqrt{(E+mc^2)/c}$$
~\\[5pt]
    
\noindent
{\bf Problem 2:} Show that:
$$\bar{u}^{(1)} u^{(1)} = 2mc$$
You don't have to calculate it, but also:
$$\bar{u}^{(2)} u^{(2)} = 2mc$$
and:
$$\bar{v}^{(1)} v^{(1)} = \bar{v}^{(2)} v^{(2)} = -2mc$$
for the anti-particle bispinors $v$.\\[5pt]

\newpage

\noindent
{\bf Problem 3:} In lecture we showed that bispinors are complete in the sense that:
$$\sum_i u^{(i)}\bar{u}^{(i)} = u^{(1)}\bar{u}^{(1)} + u^{(2)}\bar{u}^{(2)} = \slashed{p} + mc$$
and
$$\sum_i v^{(i)}\bar{v}^{(i)} = v^{(1)}\bar{v}^{(1)} + v^{(2)}\bar{v}^{(2)} = \slashed{p} - mc$$
For the simple case of particles at rest ($\vec{p} = 0$) show that this identity holds.  Hints:
In this case $p^0 = mc$ and note that:
$$\gamma^0 + 1 = \begin{pmatrix}
  2 & 0 & 0 & 0 \\
  0 & 2 & 0 & 0 \\
  0 & 0 & 0 & 0 \\
  0 & 0 & 0 & 0 \\
\end{pmatrix}$$
~\\[5pt]

\begin{figure}[htbp]
\begin{center}
\begin{feynman}
    \fermion{1.00, 1.00}{0.0, 0.2}
    \fermion{0.0, 1.8}{1.00, 1.00}
    \electroweak{1.00, 1.00}{2.00, 1.00}
    \fermion{3.00, 0.2}{2.00, 1.00}
    \fermion{2.00, 1.00}{3.00, 1.8}
    \node at (0,0) {$\bar{v}(b)$};
    \node at (0,2) {$u(a)$};
    \node at (3,2) {$\bar{u}(1)$};
    \node at (3,0) {$v(2)$};
    \node at (1.5,1.2) {$\gamma$};
    \node at (0,0.4) {$e$};
    \node at (0,1.6) {$e$};
    \node at (3.0,0.4) {$\mu$};
    \node at (3.0,1.6) {$\mu$};
\end{feynman}
\caption{\label{fig:feyn} Leading order Feynman diagram for the process $e^- + e^+ \to \gamma \to \mu^- + \mu^+$.  The particle plane-wave bi-spinor is $u$ and the anti-particle plane-wave bi-spinor is $v$.}
\end{center}
\end{figure}
~\\[5pt]

\noindent
{\bf Problem 4:} Using the Feynman rules for QED, show that the amplitude for the leading order contribution (Fig.~\ref{fig:feyn}) to $e^- + e^+ \to \mu^- + \mu^+$ is:
\begin{eqnarray*}
\mathcal{M}  &=&  -\frac{g_e^2}{(p_a+p_b)^2} \; \overline{v}(b) \, \gamma^\mu
  \, u(a)  \; \bar{u}(1) \, \gamma_\mu \, v(2) \; 
\end{eqnarray*}
where:
$$\gamma_\mu \equiv g_{\mu\nu} \gamma^\nu$$

\newpage


\noindent
{\bf Problem 5:} {\bf Casimir's Trick} In many experiments, the
incident beams are unpolarized (equal parts spin up and spin down).
Furthermore, the detectors are often incapable of determining the spin
of the final state particles.  In such contexts, the cross section is
calculated from an average over incident spins ($s_a$ and $s_b$) with four possible combinations and summed over all final state spins ($s_1$ and $s_2$).  That is, we calculate
$$ \braket{|\mathcal{M}|^2} \equiv \frac{1}{4} \; \sum_{s_a,s_b,s_1,s_2}
\; \left| \mathcal{M}(s_a,s_b,s_1,s_2)\right|^2$$
Casimir's trick yields the rather exceptionally clever and tedium avoiding identity:
$$\sum_{i,j} \;
\left( \bar{u}^{(i)}(p_a) \, \Gamma_1 \, u^{(j)}(p_b) \right) \; 
\left( \bar{u}^{(i)}(p_a) \, \Gamma_2 \, u^{(j)}(p_b) \right)^\dagger \;
= {\rm Tr}\,(\Gamma_1 \, (\slashed{p}_b + m_b c) \, \overline{\Gamma_2} \, (\slashed{p}_a + m_a c) )$$
where $\rm Tr$ indicates taking the trace, $\Gamma_1$ and $\Gamma_2$ are 4x4 matrices, and
$$\overline{\Gamma} \equiv \gamma^0 \Gamma^\dagger \gamma^0$$
Prove the identity above.  Hints:
\begin{itemize}
\item Any number is equivalently a 1x1 matrix, and so any number $x$ is identical to its trace:
$$x = {\rm Tr}(x)$$
\item The quantity
$$\left( \bar{u}^{(i)}(p_a) \, \Gamma_1 \, u^{(j)}(p_b) \right) \; 
\left( \bar{u}^{(i)}(p_a) \, \Gamma_2 \, u^{(j)}(p_b) \right)^\dagger \;$$
is just a number, and so equal to it's trace.
\item As part of your solution, you will need to show that:
$$[\bar{u}^{(i)}(p_a) \, \Gamma \, u^{(j)}(p_b)]^\dagger
= \bar{u}^{(j)}(p_b) \bar{\Gamma} u^{(i)}(p_a)$$
\item The trace is invariant under cyclic rotations, e.g.:
$$Tr(ABC) = Tr(BCA)$$
  for matrices $A$,$B$, and $C$.
\item Use the completeness of the bispinors (see previous problems).
\end{itemize}  

Note that a similar results holds if some or all of the particle spinors are replaced with anti-particle spinors.  For each antiparticle:
$$\slashed{p}+mc \to \slashed{p}-mc$$


\newpage 

\noindent
{\bf Problem 6:} For the leading order process of $e^+ + e^- \to \mu^+ + \mu^- $, and taking $m_e = m_u = 0$ show that initial state spin-averaged, final state spin summed matrix element is:
$$\braket{|\mathcal{M}|^2} = \frac{g_e^4}{4(p_a+p_b)^4}
\Tr(\gamma^\mu\slashed{p_a}\gamma^\nu\slashed{p_b})
\Tr(\gamma_\mu\slashed{p_2}\gamma_\nu\slashed{p_1})$$
Hint:  Use Casimir's trick and note that the gamma matrices have the property that:
$$\overline{\gamma_u} = \gamma_u$$
~\\[5pt]

\noindent
{\bf Problem 7:} By calculating the traces in the previous problem, show that:
$$ \braket{|\mathcal{M}|^2} 
= \frac{8 g_e^4}{(p_a+p_b)^4}
\left((p_a\cdot p_1)(p_b\cdot p_2) + (p_a\cdot p_2)(p_b\cdot p_1)\right)
$$
Hint: Griffith's works through a very similar product of traces in Example 7.6 (just note that we are taking $m$ and $M$ as zero).  You will need the identities:
$$\Tr(\gamma^\mu \gamma^\alpha \gamma^\nu \gamma^\beta) 
= 4 \, 
(    \, g^{\mu\alpha}\,g^{\nu \beta}
\, - \, g^{\mu\nu}\,g^{\alpha \beta}
\, + \, g^{\mu\beta}\,g^{\alpha\nu} \, )$$
and:
$$g_{uw}g^{uw} = 4$$

\noindent
{\bf Problem 8:} Show that in the CMS, we have:
$$ \braket{|\mathcal{M}|^2} = g_e^4 (1 + \cos^2 \theta)$$
where $\theta$ is the angle between $\vec{p}_1$ and $\vec{p}_a$. 
Hint: you want to show, for example, that:
$$p_a\cdot p_1 = \frac{E^2}{c^2}(1 - \cos \theta)$$
where $E$ is the energy of particle 1.\\[5pt]

\newpage 

\noindent
{\bf Problem 9:} Fermi's Golden Rule for $a+b \to 1+2$ scattering in the CMS and averaging over spins is:
$$ \frac{d\sigma}{d\Omega} = \left( \frac{\hbar c}{8 \pi}\right)^2 \frac{S \braket{|\mathcal{M}|^2}
}{(E_a+E_b)^2}\frac{|\vec{p}_1|}{|\vec{p}_a|}$$
Show that:
$$\frac{d\sigma}{d\Omega} = (\hbar c)^2 \; \frac{\alpha^2}{4s} \; (1 + \cos^2 \theta)$$
where $g_e = \sqrt{4\pi\alpha}$ and $s = (p_a + p_b)^2$.)

\end{document}

